\documentclass[10pt,a4paper]{amsart}
\usepackage[margin=19mm]{geometry}
\usepackage{amsmath,amssymb,mathtools}
\usepackage[scr=rsfs,cal=euler]{mathalfa}
\usepackage{xfrac}
\usepackage[unicode,hidelinks,bookmarks=false]{hyperref}
\newcommand{\ie}{i.e.\ }
\newcommand{\cf}{cf.\ }
\newcommand{\resp}{respectively\ }
\newcommand{\Subsection}{\subsection}
\providecommand{\nobreakdash}{\nobreak\hbox{-}}
\setlength{\emergencystretch}{2em}
\multlinegap0pt
\usepackage{bm}
\usepackage{bbm}
\usepackage{dsfont}
\usepackage{xspace}
\usepackage{cancel}
\usepackage{mathtools}
\usepackage{amsthm}
\makeatletter
\@ifundefined{theorem}{\newtheorem{theorem}{Theorem}}{}
\makeatother
\newcommand{\C}{\mathbb C}
\newcommand{\N}{\mathbb N}
\renewcommand{\geq}{\geqslant}
\renewcommand{\leq}{\leqslant}
\title[On low-dimensional uniform rectifiability in Heisenberg grou]{On low-dimensional uniform rectifiability in Heisenberg groups}
\author{Katrin F\"assler, Andrea Pinamonti, Kilian Zambanini}
\date{Source: arXiv:2601.03837v1 (CC BY 4.0)}
\begin{document}
\maketitle

\noindent\emph{Source and licence.} On low-dimensional uniform rectifiability in Heisenberg groups. Version arXiv:2601.03837v1 (CC BY 4.0), distributed under the Creative Commons Attribution 4.0 International licence, as recorded in the arXiv licence metadata for this exact version and shown on the abstract page https://arxiv.org/abs/2601.03837v1. Full rights notice: licenses/arxiv-2601.03837-CC-BY-4.0.txt.\par\smallskip

\noindent\textit{Editorial scope.} Counted unit: the Theorem whose source label is \texttt{injective} in the article (source0.tex lines 1371--1373, source label \texttt{injective}), delivered together with the article's own complete original proof of it (source0.tex lines 1374--1419, one \texttt{proof} environment of 46 lines). No mathematics is written, repaired or re-derived: statement and proof are the article's own text line for line. Required context retained uncounted: eq:** (lines 1355--1357); eq:* (lines 1360--1362). Every definition, display and label of those lines is retained unchanged. Omitted from this package and not redistributed: 1--1354, 1358--1359, 1363--1370, 1420--2026. Nothing inside a retained excerpt is omitted or altered: every retained body is delivered exactly as the article's lines are written, so no omission receipt of any kind accompanies this package. Citations: the retained excerpts carry no citation command at all, so no bibliography entry of the article is transcribed and no reference is redistributed. Reference adaptation: none was needed and none was made; every reference inside the delivered unit resolves inside it, so no unresolved reference or citation is produced. Printed slips kept literal: no printed word, symbol, hypothesis or number of the counted statement or of its proof was altered. Layout and class: the article's own class is \texttt{amsart} with packages including color, amsfonts, amssymb, amsmath, amscd, amstext, esint, hyperref, inputenc, graphicx, tikz, changes, comment, cleveref; the wrapper is the lane's pinned \texttt{amsart} editorial wrapper at 10pt on A4, which restates the theorem-like environments the retained text needs and adds the article's own macro definitions that the retained text needs (\texttt{\textbackslash C}, \texttt{\textbackslash N}, \texttt{\textbackslash geq}, \texttt{\textbackslash leq}), with extra wrapper packages (bm, bbm, dsfont, xspace, cancel, mathtools). The delivered numbering is the unit's own. Assets: no retained span contains a figure environment, a graphics-inclusion command, a PDF-page inclusion, a table or a TikZ picture; no image file is copied into this package, the package declares no asset dependency and no image file is redistributed with it. Licence: version arXiv:2601.03837v1 of this article is distributed under the Creative Commons Attribution 4.0 International licence, as recorded in the arXiv licence metadata for this exact version and shown on the abstract page https://arxiv.org/abs/2601.03837v1. A full rights notice is delivered at licenses/arxiv-2601.03837-CC-BY-4.0.txt. Characters: every retained excerpt is pure ASCII, so the delivered engine, XeLaTeX with its default Latin Modern text font, reports no missing character.
\begin{equation}\label{eq:**}
    2\cdot4^{-n}\leq l_n\leq 2^{-n+1}\text{ for every }n\geq 0.
\end{equation}

\begin{equation}\label{eq:*}
    L \leq 2 \cdot \exp(0.08)\leq 2.2.
\end{equation}

\begin{theorem}\label{injective}
    The curve $\omega^\C$ (hence also  $\omega$) is injective, and $\Gamma$ is 1-regular in $(\mathbb{H}^1,d)$.
\end{theorem}

\begin{proof}
The proof of injectivity relies on the following claim:\medskip\\
  \textbf{Claim}. Let $p=\omega^\C(\sigma/4^n)$ for some $\sigma=0,1,\dots, 4^n$. Then
  \begin{equation}\label{eq:claim} B_{\rm Eucl}(p,0.6\, l_n)\,\cap\,\omega^\C\left([0,1]\setminus\left[\frac{\sigma-1}{4^n},\frac{\sigma +1}{4^n}\right]\right)=\emptyset. 
  \end{equation}
We now assume the claim and prove that $\omega^\C$ is injective. Let $0\leq a<b\leq 1$ and choose $n\in\N$, $\sigma=0,1,\dots, 4^n$ such that 
\begin{equation}\label{distancefroma}
    \left|a-\frac{\sigma}{4^n}\right|\leq \frac{1}{2\cdot4^n}
\end{equation}
and 
\begin{equation}\label{distancefromb}
   \left|b-\frac{\sigma}{4^n}\right|> \frac{1}{4^n}. 
\end{equation}
This is always possible, since one can choose $n$ large enough such that
\[|a-b|> \frac{3}{2\cdot4^n}\] and subsequently choose $\sigma=0,1,\dots,4^n$ such that \eqref{distancefroma} holds. By the triangular inequality, \eqref{distancefromb} holds as well.
From \eqref{distancefroma} it follows that 
\begin{equation}\label{eq:ainball}\left|\omega^\C(a)-\omega^\C\left(\frac{\sigma}{4^n}\right)\right|\leq L\left|a-\frac{\sigma}{4^n}\right|\leq L\,\frac{1}{2\cdot 4^n}\overset{\eqref{eq:*}}{\leq} \frac{2.4}{2\cdot 4^n}=0.6\frac{2}{4^n}\overset{\eqref{eq:**}}{\leq} 0.6 \,l_n.\end{equation} Hence $\omega^\C(a)\in B_{\rm Eucl}(p,0.6\, l_n)$. On the other hand,  \eqref{distancefromb} means that $b\in[0,1]\setminus [\frac{\sigma-1}{4^n},\frac{\sigma +1}{4^n}]$. Therefore, the claim implies that $\omega^\C(b)\not\in B_{\rm Eucl}(p,0.6\, l_n)$ and in particular we deduce $\omega^\C(a)\neq \omega^\C(b)$, proving injectivity of $\omega^\C$ (and clearly of $\omega$).\medskip\\
 Assuming the claim, we now show that $\omega([0,1])$ is 1-regular. Fix now $x=(z,t)\in \omega([0,1])$ and $r>0$. In order to prove upper 1-regularity, we can assume that $r<1/20$, since otherwise we use the simple estimate
  \[\mathcal H^1\left(B(x,r)\,\cap\,\omega\left([0,1]\right)\right)\leq  L\leq 20 L\,r.\]
For $r<1/20$,  choose $n\in\N$ such that $l_{n+1}/40\leq r<l_n/40$.
  By $L$-Lipschitz continuity of $\omega^\C$, there exists $\sigma=0,1,\dots 4^n$ such that $|z-p|\leq \frac{L}{2\cdot 4^n}$, where $p=\omega^\C(\sigma/4^n)$. In particular, %\textcolor{blue}{(Here explain why we can assume $L\leq 2.2$ instead of 2.4. Maybe already in the construction of the curve)}
  \[|z-p|\leq \frac{L}{2\cdot 4^n}\overset{\eqref{eq:*}}{\leq} \frac{2.2}{2\cdot 4^n}=\frac{11}{20}\frac{2}{4^n}\overset{\eqref{eq:**}}{\leq} \frac{11}{20}l_n.\]
  Hence, if $z'\in B_{\rm Eucl}(z,r)$, then
  \[|z'-p|\leq |z'-z|+|z-p|\leq r+\tfrac{11}{20} l_n\leq \left(\tfrac{1}{40}+\tfrac{11}{20}\right)l_n\leq 0.6\, l_n.\]
  This proves that $B_{\rm Eucl}(z,r)\subseteq B_{\rm Eucl}(p,0.6\,l_n)$. Moreover, \eqref{eq:claim} implies that
  \begin{equation}\label{eq:BallCurveIncl}
B_{\rm Eucl}(p,0.6\, l_n)\,\cap\,\omega^\C\left([0,1]\right)\subseteq \omega^\C\left(\left[\frac{\sigma-1}{4^n},\frac{\sigma +1}{4^n}\right]\right).
  \end{equation}
  Altogether, this shows that
  \begin{align*}
      \pi\left(B(x,r)\right)\cap \pi\left(\omega([0,1])\right)
     & = \pi\left(B(x,r)\right)\cap \omega^{\mathbb{C}}\left([0,1] \right)=B_{\mathrm{Eucl}}(z,r)\cap \omega^{\mathbb{C}}\left([0,1]\right)\\
  &{=}  B_{\rm Eucl}(p,0.6\,l_n)\cap \omega^{\mathbb{C}}\left([0,1]\right)
 \overset{\eqref{eq:BallCurveIncl}}{=} \omega^\C\left(\left[\frac{\sigma-1}{4^n},\frac{\sigma +1}{4^n}\right]\right).
  \end{align*}
Therefore, $\pi\left(B(x,r)\cap \omega([0,1])\right)\subset \omega^\C\left(\left[\frac{\sigma-1}{4^n},\frac{\sigma +1}{4^n}\right]\right)$. The injectivity of $\omega^{\mathbb{C}}$  implies then that also
\begin{equation}\label{eq:CovByCurv}
    B(x,r)\cap \omega([0,1]) \subset \omega\left(\left[\frac{\sigma-1}{4^n},\frac{\sigma +1}{4^n}\right]\right).
\end{equation}
Moreover,
by the $L$-Lipschitz continuity of $\omega^\C$ and the same estimates as in \eqref{eq:ainball}, we have \begin{equation}\label{1regularity}\mathcal{H}^1_{\mathrm{Eucl}}\left(\omega^\C\left(\left[\frac{\sigma-1}{4^n},\frac{\sigma +1}{4^n}\right]\right)\right)\leq L\,\frac{2}{4^n}\overset{\eqref{eq:**},\eqref{eq:*}}{\leq} 2.4\,l_n.\end{equation}
The upper 1-regularity of $\omega([0,1])$ then follows since, by \eqref{eq:CovByCurv} and the horizontality of the curve $\omega$,
  \[\begin{split}\mathcal H^1(B(x,r)\cap\, \omega([0,1]))\leq \mathcal H^1\left(\omega\left(\left[\frac{\sigma-1}{4^n},\frac{\sigma +1}{4^n}\right]\right)\right)&= \mathcal{H}^1_{\mathrm{Eucl}}\left(\omega^{\mathbb{C}}\left(\left[\frac{\sigma-1}{4^n},\frac{\sigma +1}{4^n}\right]\right)\right)\\&\overset{\eqref{1regularity}}{\leq} 2.4\cdot 4\cdot 40 \,r.\end{split}\]
 
  The lower regularity of $\omega([0,1])$ is immediate since it is a  connected set.
\end{proof}

\end{document}
